% TOP_PROBLEM: {"id":30007522,"problem_number":"LOCAL-30007522","title":"Supply-versus-demand inflation","category_id":21,"category":{"id":21,"name":"economics","display_name":"Economics","description":"Open economic theory statements, published research agendas, and proposed scoped specifications; question provenance and historical priority are qualified per record.","slug":"economics","order_index":21,"created_at":"2026-10-08T00:00:00Z"},"status":"provenance_pending","external_url":null,"legacy_ids":[],"published":false,"statement":"Identify convention-dependent supply, demand, and reallocation contributions to inflation, allocating nonlinear interactions by a stated Shapley rule.","economics":{"original_id":11,"field":"Inflation and monetary policy","kind":"identification","formalization_basis":"proposed_specification","source_status":"not_verified","priority_status":"not_established","priority_note":"No explicit primary open-problem statement matching this catalogue item was verified; supporting research is not evidence of first formulation.","earliest_verified_reference":null,"current_reference":{"authors":["Julian di Giovanni","Şebnem Kalemli-Özcan","Alvaro Silva","Muhammed A. Yıldırım"],"title":"Pandemic-Era Inflation Drivers and Global Spillovers","year":2023,"url":"https://www.newyorkfed.org/medialibrary/media/research/staff_reports/sr1080.pdf","doi":"10.59576/sr.1080","locator":"Abstract and Section 1, pp. 1–3; July 2024 revision","evidence_summary":"The paper quantifies pandemic inflation and international spillovers in a specified multicountry production-network model. It supplies a partial answer and framework, not an explicit general open-problem statement about business-cycle synchronisation or inflation decomposition."},"formalization_note":"New decomposition specification. The cited paper already provides a pandemic-specific model-based answer; it is not an original open-problem statement."},"statusQualification":"Provenance pending: an explicit published statement of this exact open question was not verified. This is a catalogue-authored candidate specification, not a certified open conjecture. New decomposition specification. The cited paper already provides a pandemic-specific model-based answer; it is not an original open-problem statement.","statusReviewedAt":"2026-10-08"}
% FIRST_POSTED: 2026-10-08
% ENTRY_KIND: statement_only
% !TeX program = lualatex
\documentclass[11pt]{article}
\usepackage{fontspec}
\usepackage[a4paper,margin=25mm]{geometry}
\usepackage{amsmath,amssymb,mathtools}
\usepackage{xurl}
\usepackage[hidelinks]{hyperref}
\newcommand{\sref}[2]{\hyperlink{source:#1:#2}{[S#2]}}
\setlength{\emergencystretch}{3em}
\begin{document}
% problemId:30007522
\section*{Supply-versus-demand inflation}\label{30007522}
\subsection{Definitions and mathematical statement}
Fix a multicountry, multisector economy with observed sectoral prices, wages, output, expenditures, and input–output links. Its structural model maps a finite sequence of shocks \(e\) and parameter \(\theta\) into aggregate inflation \(\pi_t(e;\theta)\). Partition shocks into supply, aggregate demand, and sectoral demand reallocation. State the instruments or exclusion restrictions distinguishing these classes; price–quantity comovement alone must not be assumed to establish causal classification. Let \(e^S\) retain only shocks in subset \(S\) of the three classes, with a common initial state and policy rule.
To allocate nonlinear interactions, define a Shapley contribution for class \(k\):
\[
C_{k,t}=\sum_{S\subseteq K\setminus\{k\}}
\frac{|S|!(|K|-|S|-1)!}{|K|!}
\big[\pi_t(e^{S\cup\{k\}};\theta)-\pi_t(e^S;\theta)\big],
\]
where \(K\) is the three-class set. Identify or bound these contributions over a declared inflation episode, including uncertainty in shocks, elasticities, and network links. Test whether decompositions remain stable under plausible alternative identification schemes and inflation aggregation weights. Counterfactual policy responses must be declared as fixed rules or endogenous responses. The target is an explicitly convention-dependent causal decomposition, not a unique model-free attribution of each observed price increase to supply or demand.

\par\textbf{Formulation and scope.} New decomposition specification. The cited paper already provides a pandemic-specific model-based answer; it is not an original open-problem statement.

\par\textbf{Open-question provenance.} An explicit published open-question statement was not verified. Sources below are supporting literature and must not be treated as a first listing.

\par\textbf{Historical priority.} No explicit primary open-problem statement matching this catalogue item was verified; supporting research is not evidence of first formulation.

\subsection{Short English statement}
Identify convention-dependent supply, demand, and reallocation contributions to inflation, allocating nonlinear interactions by a stated Shapley rule.

\subsection{Sources}
\begin{itemize}\raggedright
\item[\textnormal{[S1]}]\hypertarget{source:30007522:1}{} Julian di Giovanni, Şebnem Kalemli-Özcan, Alvaro Silva, Muhammed A. Yıldırım. 2023. Pandemic-Era Inflation Drivers and Global Spillovers. Abstract and Section 1, pp. 1–3; July 2024 revision.
\url{https://www.newyorkfed.org/medialibrary/media/research/staff_reports/sr1080.pdf}.
DOI: 10.59576/sr.1080.
{\small\textit{Supporting or current literature; see evidence qualification. The paper quantifies pandemic inflation and international spillovers in a specified multicountry production-network model. It supplies a partial answer and framework, not an explicit general open-problem statement about business-cycle synchronisation or inflation decomposition.}}
\end{itemize}
\end{document}
